\documentclass{report}
\usepackage[spanish]{babel}
\usepackage[utf8]{inputenc}
\usepackage{graphicx, longtable, float, titlesec, hyperref, enumitem, dingbat, soul, multicol, listings}
\usepackage[dvipsnames]{xcolor}
\usepackage[margin=2cm]{geometry}

% Cambia el color de los links
\hypersetup{hidelinks}

% Generamos un comando para saltar pagina con las secciones
\NewDocumentCommand{\cpsection}{s o m}{%
  \clearpage
  \IfBooleanTF{#1}
    {\section*{#3}}
    {%
      \IfNoValueTF{#2}
        {\section{#3}}
        {\section[#2]{#3}}%
    }%
}

% Python Code
\lstdefinestyle{Python}{
  commentstyle=\color{brown},
  keywordstyle=\color{violet},
  numberstyle=\tiny\color{gray},
  stringstyle=\color{purple},
  basicstyle=\ttfamily\footnotesize,
  breakatwhitespace=false,         
  breaklines=true,                 
  captionpos=b,                    
  keepspaces=true,                 
  numbers=left,                    
  numbersep=5pt,                  
  showspaces=false,                
  showstringspaces=false,
  showtabs=false,                  
  tabsize=2,
  literate={ñ}{{\~n}}1 {á}{{\'a}}1 {é}{{\'e}}1 {í}{{\'i}}1 {ó}{{\'o}}1 {ú}{{\'u}}1
}
\lstset{style=Python}

% Elimina la palabra "Capítulo" de los títulos de los capítulos
\titleformat{\chapter}[display]
  {\normalfont\bfseries}{}{0pt}{\Huge\thechapter.\space}

\titleformat{name=\chapter,numberless}[display]
  {\normalfont\bfseries}{}{0pt}{\Huge}

\titlespacing*{\chapter}{0pt}{-50pt}{20pt}

% Personalización del índice de listados
\renewcommand{\lstlistingname}{Código}  % Cambiar el nombre de "Listing" a "Código"
\renewcommand{\lstlistlistingname}{Índice de Códigos}

% Añade numeración a los subsubsection*s y los añade al índice
\setcounter{secnumdepth}{4}
\setcounter{tocdepth}{4}

\begin{document}
  \begin{titlepage}
      \centering
      \includegraphics[width=0.6\textwidth]{./.img/logo.jpg}\\
      \vspace{1cm}
      \LARGE Técnicas de Inteligencia Artificial\\
      \vspace{0.5cm}
      \Large Ingeniería Informática de Gestión y Sistemas de Información\\
      \vspace{3cm}
      \Huge Practica 3\\
      \huge Clasificacion\\
      \vspace{2.5cm}
      \Large Autor(es):\\
      \vspace{0.2cm}
      \large Xabier Gabiña\\
      \large Diego Montoya\\
      \vfill
      \today
  \end{titlepage}
  \tableofcontents
  %\listoffigures
  %\listoftables
  \lstlistoflistings
  \chapter{Introducción}
    \paragraph*{}{
      En este proyecto, se abordará la implementación de algoritmos de clasificación utilizando el modelo de perceptrón, enfocado en resolver problemas de lógica y clonación de comportamiento en un entorno de aprendizaje. El desarrollo comienza con la construcción de puertas lógicas básicas (AND, OR, NOT y XOR) mediante perceptrones para comprender el funcionamiento del modelo en conjuntos de datos limitados. Posteriormente, se extiende su aplicación hacia la creación de clasificadores más complejos capaces de reconocer dígitos y replicar comportamientos observados en un agente de juego.
    }
  \chapter{Ejercicios}
    \section{Puertas Lógicas con perceptron} %TERMINADO
      \subsection*{Descripción}
        \paragraph*{}{
          En este apartado se implementaran las puertas lógicas AND, OR, NOT y XOR mediante perceptrones.
          Empezaremos haciendo pruebas ya conociendo los pesos para pasar a entrenar nosotros mismos los pesos y para finalizar utilizaremos el weighted average.
        }
      \subsection*{Implementación}
        \begin{lstlisting}[language=Python, caption=Implementación del perceptron para las puerta lógica AND]
# Definir dos vectores (listas): input my_x, pesos my_w
my_x = [0, 1]#input un item
my_w = [0.66, 0.80]

# Multiplicar dos vectores elemento a elemento
def mul(a, b):
    """
    Devolver una lista c, de la misma longitud que a y b, donde 
    cada elemento c[i] = a[i] * b[i]
    """
    return [a[i] * b[i] for i in range(len(a))]

mul(my_x, my_w)

my_bias  = 1
my_wbias = -0.97

my_wPlusWBias = [my_wbias] + my_w

def distanciaDelCoseno(x, weights, bias):
    """
    El producto escalar (producto punto) de dos vectores y la similitud de coseno no son completamente equivalentes 
    ya que la similitud del coseno solo se preocupa por la diferencia de ángulo, 
    mientras que el producto de punto se preocupa por el ángulo y la magnitud
    Pero en muchas ocasiones se emplean indistintamente
    Así pues, esta función devuelve el valor escalar de la neurona, es decir, 
    el producto escalar entre el vector de entrada añadiendo el bias y el vector de los pesos
    recordad que "sum(list)" computa la suma de los elementos de una lista
    Así pues se comenzará por añadir el bías en la posición 0 del vector de entrada 
    antes de llevar a cabo el producto escalar para así tener dos vectores de 
    la misma longitud. Emplea la función mul que ya has programado
    """
    x = [bias] + x
    return sum(mul(x, weights))

distanciaDelCoseno(my_x, my_wPlusWBias, my_bias)

def neuron(x, weights, bias):
    """
    Devolverá el output de una neurona clásica 
    (reutilizar la distancia del coseno definida previamente) 
    y añadir la función de activación (step function): si >=0 entonces 1 sino -1
    """
    return 1 if distanciaDelCoseno(x, weights, bias)>=0 else -1

neuron(my_x, my_wPlusWBias, my_bias)

def and_neuron(x):
    """
    Devuelve x1 AND x2 suponiendo que la hemos entrenado
    y que en ese entrenamiento hemos aprendido los pesos apropiados 
    (mirar las transparencias de clase). Así pues inicializaremos 
    una la variable local and_w con los pesos aprendidos 
    y a 1 la variable local and_bias 
    y ejecutaremos la función neurona para el item x"""
    and_w    = [-0.97,0.66, 0.80]#initialization of the weights and_w
    and_bias = 1#initialization of the bias and_bias
    return neuron(x, and_w, and_bias)

my_x_collection = [
    [0, 0],
    [0, 1],
    [1, 0],
    [1, 1],
]

print('Testando el output de la neurona AND')
#bucle para ir obteniendo el output de la neurona AND para cada item del input
for my_x in my_x_collection:
    print(my_x, f'{and_neuron(my_x):.3f}')
        \end{lstlisting}
        \clearpage
        \begin{lstlisting}[language=Python, caption=Implementación del perceptron para las puerta lógica OR]
from random import seed, random

# Inicialización
print('Entrenando una neurona OR hasta convergencia')
notConverge=True
seed(1)

orWeights= [random() for i in range(3)]
orBias = 1
orGoldOutputs = [-1,1,1,1]

# Entrenamiento
numeroVuelta = 0
while notConverge:
    notConverge = False
    for i, my_x in enumerate(my_x_collection):
        x_with_bias = [orBias] + my_x  # Agrega el bias al vector de entrada
        predicted_output = neuron(my_x, orWeights, orBias)
        
        # Si la predicción no coincide con la salida esperada, ajustar pesos
        if predicted_output != orGoldOutputs[i]:
            adjustment = 1 if orGoldOutputs[i] == 1 else -1
            orWeights = [orWeights[j] + adjustment * x_with_bias[j] for j in range(len(orWeights))]
            notConverge = True  # Continuar iterando hasta que todo esté correcto
    numeroVuelta += 1
    print(f"Vuelta {numeroVuelta}: Pesos actualizados OR:", orWeights)
        \end{lstlisting}
        \clearpage
        \begin{lstlisting}[language=Python, caption=Implementación del perceptron para las puerta lógica NOT]
my_x_collection = [
    [0],
    [1]
]
from random import seed, random


# Inicializaciones
print('Entrenando una neurona NOT hasta convergencia')
notConverge = True
seed(1)

notWeights = [random(), random()]
notBias   = 1
notGoldOutput = [1, -1]

#entrenando
numeroVuelta = 0
while notConverge:
    notConverge = False
    for i, my_x in enumerate(my_x_collection):
        x_with_bias = [orBias] + my_x  # Agrega el bias al vector de entrada
        predicted_output = neuron(my_x, orWeights, orBias)
        
        # Si la predicción no coincide con la salida esperada, ajustar pesos
        if predicted_output != orGoldOutputs[i]:
            adjustment = 1 if orGoldOutputs[i] == 1 else -1
            orWeights = [orWeights[j] + adjustment * x_with_bias[j] for j in range(len(orWeights))]
            notConverge = True  # Continuar iterando hasta que todo esté correcto
    numeroVuelta += 1
    print(f"Vuelta {numeroVuelta}: Pesos actualizados NOT:", orWeights)

        \end{lstlisting}
        \clearpage
        \begin{lstlisting}[language=Python, caption=Implementación del perceptron para las puerta lógica OR con Weighted average]
my_x_collection = [
    [0, 0],
    [0, 1],
    [1, 0],
    [1, 1],
]

def matrixAverage(m):
  res=list()
  acum=list()
  if len(m) > 0:
      res=[0]*len(m[0])
      for v in m:
          res = [a+b for a,b in zip (res,v)]
      acum=[elem/len(m) for elem in res]
  return acum

matrix=[[2,3,4],[2,3,4],[2,3,4]]
print(matrixAverage(matrix))

from random import seed, random

# Inicializaciones
print('Entrenando una neurona OR hasta convergencia')
notConverge=True
seed(1)

orWeights= [random() for i in range(3)]
orBias = 1
orGoldOutputs = [-1,1,1,1]

# Historial de pesos
weightHistory = []

# Entrenamiento
numeroVuelta = 0
while notConverge:
    notConverge = False
    epoch_weights = orWeights.copy()  # Copia de los pesos para acumular en cada epoch
    for i, my_x in enumerate(my_x_collection):
        x_with_bias = [orBias] + my_x  # Agrega el bias al vector de entrada
        
        # Verificar si la predicción coincide con la salida esperada
        if neuron(my_x, orWeights, orBias) != orGoldOutputs[i]:
            # Ajuste de pesos basado en el error
            adjustment = 1 if orGoldOutputs[i] == 1 else -1
            orWeights = [orWeights[j] + adjustment * x_with_bias[j] for j in range(len(orWeights))]
            notConverge = True  # Continuar iterando hasta convergencia
    weightHistory.append(orWeights.copy())  # Guardar los pesos de esta epoch
    numeroVuelta += 1
    print(f"Vuelta {numeroVuelta}: Pesos actualizados OR:", orWeights)

# Calcular el promedio ponderado de los pesos al finalizar el entrenamiento
averageWeights = matrixAverage(weightHistory)
print("\nPromedio ponderado de los pesos OR:", averageWeights)

def or_neuron(x):
    """
    Devuelve x1 AND x2 suponiendo que la hemos entrenado
    y que en ese entrenamiento hemos aprendido los pesos apropiados 
    (mirar las transparencias de clase). Así pues inicializaremos 
    una la variable local and_w con los pesos aprendidos 
    y a 1 la variable local and_bias 
    y ejecutaremos la función neurona para el item x"""
    or_w    = [-0.3656,0.8474, 0.7637]#initialization of the weights and_w
    or_bias = 1#initialization of the bias and_bias
    return neuron(x, or_w, or_bias)

print('Testando el output de la neurona OR')
#bucle para ir obteniendo el output de la neurona AND para cada item del input
for my_x in my_x_collection:
    print(f'Input: {my_x} -> OR Output: {or_neuron(my_x):.3f}')
        \end{lstlisting}
        \clearpage
        \begin{lstlisting}[language=Python, caption=Implementación del perceptron para las puerta lógica XOR]
# Combinando una puerta OR y una AND, y aprendiendo el peso que hay que darle a cada una para obtener un XOR 
from random import seed, random

# Inicializaciones
print('Entrenando una neurona XOR hasta convergencia')
xorConverge=True
seed(1)

xorWeights= [random() for i in range(3)]
xorBias   = -0.5
xorGoldOutputs=[1,-1,-1,1]

# Entrenando
numeroVuelta = 0
while xorConverge:
    xorConverge = False
    for i, my_x in enumerate(my_x_collection):
        # Usar las salidas de las puertas AND y OR como entradas para la XOR
        and_output = and_neuron(my_x)
        or_output = or_neuron(my_x)
        
        # Crear el nuevo vector de entrada con los resultados de AND y OR
        new_x = [xorBias, and_output, or_output]
        
        # Verificar si la predicción coincide con el valor esperado
        if neuron([and_output, or_output], xorWeights, xorBias) != xorGoldOutputs[i]:
            # Ajuste de pesos
            adjustment = 1 if xorGoldOutputs[i] == 1 else -1
            xorWeights = [xorWeights[j] + adjustment * new_x[j] for j in range(len(xorWeights))]
            xorConverge = True  # Continuar iterando hasta convergencia
    numeroVuelta += 1
    print(f"Vuelta {numeroVuelta}: Pesos actualizados XOR:", xorWeights)

def xor_neuron(x):
  """
  Return x1_ * x2 + x1 * x2_
  """
  xor_w    = [-1.115635755887599, 0.3474337369372327, -0.7362253810233859]
  xor_bias = -0.5
  new_x=list()
  new_x.append(and_neuron(x))
  new_x.append(or_neuron(x))
  return neuron(new_x, xor_w, xor_bias)

print('Checking XOR neuron output')
for my_x in my_x_collection:
    print(my_x, f'{xor_neuron(my_x):.3f}')
        \end{lstlisting}
        \clearpage
      \subsection*{Conclusiones}
        \paragraph*{}{
          Respecto a la implementación de AND ha sido sencilla ya que lo más complicado que es encontrar los pesos no hemos tenido que realizarlo por lo que no hay mucho que comentar al respecto.\\
          
          En cuanto a la de OR y NOT de a continuación, ya hemos tenido que realizar el entrenamiento de los pesos. Aqui hemos tenido algun problema con el tema de añadir el BIAS y la longitud de los vectores, pero una vez solucionado hemos podido realizar el entrenamiento de los pesos sin mayor problema.\\

          En cuanto al weighted average, es practicamente el mismo entrenamiento pero guardando los pesos en cada iteración para luego calcular la media ponderada de los pesos.\\

          Por último, la implementación de XOR ha sido la más complicada ya que hemos tenido que combinar las puertas OR y AND para obtener el resultado de la XOR. Una vez obtenido el resultado de las puertas OR y AND, hemos tenido que realizar el entrenamiento de los pesos para obtener el resultado de la XOR pero si nos ha costado ver como podiamos combinar las puertas OR y AND para obtener el resultado de la XOR.\\
        }
    \cpsection{Perceptron}
      \subsection*{Descripción}
        \paragraph*{}{
          El ejercicio consiste en implementar una función de entrenamiento para un modelo de perceptrón, con el objetivo de clasificar digitos en diferentes categorías.
          Este modelo de perceptrón es un tipo de algoritmo de aprendizaje supervisado que ajusta un vector de pesos en función de los errores de clasificación en el conjunto de entrenamiento.\\

          El perceptrón funciona de la siguiente manera:
          \begin{enumerate}
            \item Obtenemos los datos y los pesos asociados a cada etiqueta.
            \item Iteramos sobre los datos de entrenamiento y calculamos los puntajes para cada etiqueta. Este puntuaje se consigue multiplicando el vector de características por el vector de pesos de cada etiqueta.
            \item Predecimos la etiqueta con el puntaje más alto.
            \item Si la predicción es incorrecta, actualizamos los pesos de la etiqueta correcta y de la etiqueta predicha.
            \item Repetimos el proceso hasta que se alcance el número máximo de iteraciones o hasta que no haya errores de clasificación.
          \end{enumerate}
        }
      \subsection*{Implementación Final}
        \begin{lstlisting}[language=Python, caption=Implementación final del perceptron]
class PerceptronClassifier:
  """
  Perceptron classifier.

  Note that the variable 'datum' in this code refers to a counter of features
  (not to a raw samples.Datum).
  """

  def __init__(self, legalLabels, max_iterations):
      self.legalLabels = legalLabels
      self.type = "perceptron"
      self.max_iterations = max_iterations
      self.weights = {}
      self.features = None
      for label in legalLabels:
          self.weights[label] = util.Counter()  # this is the data-structure you should use

  def setWeights(self, weights):
      assert len(weights) == len(self.legalLabels)
      self.weights = weights

  def train(self, trainingData, trainingLabels, validationData, validationLabels):
      """
      The training loop for the perceptron passes through the training data several
      times and updates the weight vector for each label based on classification errors.
      See the project description for details.

      Use the provided self.weights[label] data structure so that
      the classify method works correctly. Also, recall that a
      datum is a counter from features to values for those features
      (and thus represents a vector a values).
      """

      self.features = trainingData[0].keys()  # could be useful later
      # DO NOT ZERO OUT YOUR WEIGHTS BEFORE STARTING TRAINING, OR
      # THE AUTOGRADER WILL LIKELY DEDUCT POINTS.

      for iteration in range(self.max_iterations):
          print("Starting iteration ", iteration, "...")
          for i in range(len(trainingData)):  # training data
              # Obtener el ejemplo y su etiqueta real
              x_i = trainingData[i]
              y_i = trainingLabels[i]

              # Calcular los puntajes para cada etiqueta
              scores = util.Counter()
              for label in self.legalLabels:
                  scores[label] = self.weights[label] * x_i

              # Predecir la etiqueta con el puntaje más alto
              predicted_label = scores.argMax()

              # Si la predicción es incorrecta, actualizar los pesos
              if predicted_label != y_i:
                  self.weights[y_i] += x_i
                  self.weights[predicted_label] -= x_i

  def classify(self, data):
      """
      Classifies each datum as the label that most closely matches the prototype vector
      for that label.  See the project description for details.

      Recall that a datum is a util.counter...
      """
      guesses = []
      for datum in data:
          vectors = util.Counter()
          for label in self.legalLabels:
              vectors[label] = self.weights[label] * datum
          guesses.append(vectors.argMax())
      return guesses

  def findHighWeightFeatures(self, label):
      """
      Returns a list of the 100 features with the greatest weight for some label
      """
      # Obtener los pesos de los features para la etiqueta dada
      weights = self.weights[label]

      # Ordenar los features por peso en orden descendente
      sorted_features = weights.sortedKeys()

      # Seleccionar los 100 features con mayor peso
      featuresWeights = sorted_features[:100]

      return featuresWeights
        \end{lstlisting}
      \subsection*{Conclusiones}
        \paragraph*{}{
          La implementación del perceptron es bastante diferente a lo que hemos empezado haciendo en el ejercicio de las puertas lógicas aunque la idea es la misma.
          La diferencia es que ya no tenemos una clasificación binaria sino que tenemos varias etiquetas y por lo tanto tenemos que calcular los puntajes para cada etiqueta y predecir la etiqueta con el puntaje más alto.
          Esto, tambien provoca que las actualizaciones de los pesos sean diferentes ya que ahora tenemos que actualizar los pesos de la etiqueta correcta y de la etiqueta predicha.
          De esta forma, sumando a la etiqueta correcta y restando a la etiqueta predicha conseguimos que la etiqueta correcta se vuelve cada vez más probable en comparación con las etiquetas incorrectas.
        }
    \cpsection{Clonando el Comportamiento del Pacman}
      \subsection*{Descripción}
        \paragraph*{}{
          En este ejercicio se implementará un clonador de comportamiento para el agente Pacman, utilizando un modelo de perceptrón para clasificar los movimientos legales en un estado de juego.
          El objetivo es entrenar el clonador para que pueda predecir el movimiento correcto del Pacman en función de las características de un estado de juego.\\
        }
      \subsection*{Primera Implementación}
        \begin{lstlisting}[language=Python, caption=Implementación inicial del clonador de comportamiento del pacman]
class PerceptronClassifierPacman(PerceptronClassifier):
    def __init__(self, legalLabels, maxIterations):
        PerceptronClassifier.__init__(self, legalLabels, maxIterations)
        self.weights = util.Counter()

    def classify(self, data):
        """
        Data contains a list of (datum, legal moves)
        
        Datum is a Counter representing the features of each GameState.
        legalMoves is a list of legal moves for that GameState.
        """
        guesses = []
        for datum, legalMoves in data:
            vectors = util.Counter()
            for move in legalMoves:
                vectors[move] = self.weights * datum[move]  # changed from datum to datum[l]
            guesses.append(vectors.argMax())
        return guesses

    def train(self, trainingData, trainingLabels, validationData, validationLabels):
        self.features = trainingData[0][0]['Stop'].keys()  # could be useful later
        # DO NOT ZERO OUT YOUR WEIGHTS BEFORE STARTING TRAINING, OR
        # THE AUTOGRADER WILL LIKELY DEDUCT POINTS.

        for iteration in range(self.max_iterations):
            print("Starting iteration ", iteration, "...")
            for i in range(len(trainingData)):
                #print(trainingData[i])
                # Obtener el ejemplo, el movimiento legal y la etiqueta
                x_i, legalMoves = trainingData[i]
                y_i = trainingLabels[i]
                
                # Calcular los puntajes para cada etiqueta
                scores = util.Counter()
                for move in legalMoves:
                    #print(x_i[move])
                    scores[move] = self.weights[move] * x_i[move]['foodCount']
                
                # Obtener la etiqueta con el puntaje más alto
                predicted_move = scores.argMax()
                
                # Si la predicción es incorrecta, actualizar los pesos
                if y_i != predicted_move:
                    self.weights[y_i] += x_i[predicted_move]['foodCount']
                    self.weights[predicted_move] -= x_i[predicted_move]['foodCount']
        \end{lstlisting}
        \vfill
      \subsection*{Final Implementación}
        \begin{lstlisting}[language=Python, caption=Implementación final del clonador de comportamiento del pacman]
class PerceptronClassifierPacman(PerceptronClassifier):
    def __init__(self, legalLabels, maxIterations):
        PerceptronClassifier.__init__(self, legalLabels, maxIterations)
        self.weights = util.Counter()

    def classify(self, data):
        """
        Data contains a list of (datum, legal moves)
        
        Datum is a Counter representing the features of each GameState.
        legalMoves is a list of legal moves for that GameState.
        """
        guesses = []
        for datum, legalMoves in data:
            vectors = util.Counter()
            for move in legalMoves:
                vectors[move] = self.weights * datum[move]  # changed from datum to datum[l]
            guesses.append(vectors.argMax())
        return guesses

    def train(self, trainingData, trainingLabels, validationData, validationLabels):
        self.features = trainingData[0][0]['Stop'].keys()  # could be useful later
        # DO NOT ZERO OUT YOUR WEIGHTS BEFORE STARTING TRAINING, OR
        # THE AUTOGRADER WILL LIKELY DEDUCT POINTS.

        for iteration in range(self.max_iterations):
            print("Starting iteration ", iteration, "...")
            for i in range(len(trainingData)):
                # Obtener el ejemplo, el movimiento legal y la etiqueta
                x_i, legalMoves = trainingData[i]
                y_i = trainingLabels[i]
                
                # Calcular los puntajes para cada etiqueta
                scores = util.Counter()
                for move in legalMoves:
                    scores[move] = self.weights * x_i[move]
                
                # Obtener la etiqueta con el puntaje más alto
                predicted_move = scores.argMax()
                
                # Si la predicción es incorrecta, actualizar los pesos
                if y_i != predicted_move:
                    self.weights += x_i[y_i]
                    self.weights -= x_i[predicted_move]
        \end{lstlisting}
      \subsection*{Conclusiones}
        \paragraph*{}{
          En este ejercicio nos hemos hecho un poco lio con las estructuras, especialmente, con las del peso (Por no leer bien el enunciado para variar).
          Lo principal es que ahora trainingData es una lista con dos elementos, el primero es un diccionario con los movimientos y sus valores y el segundo es una lista con los movimientos legales.
          Tambien el self.weights es un vector de pesos unicos en vez de un diccionario ya que ahora estas con las clasificación de decisiones por lo que el vector de pesos es compartido mientras que antes era un diccionario con los pesos de cada etiqueta.\\

          Teniendo en cuenta esto, la implementación es bastante similiar a la del Perceptron del ejercicio anterior.
        }
    \cpsection{Clonando el Comportamientodel Pacman con rasgos diseñados por nosotros}
      \subsection*{Descripción}
        \paragraph*{}{
          En este ejercicio tendremos que definir nuestros propios rasgos para el clonador de comportamiento del Pacman y entrenar el modelo de perceptrón con estos rasgos.
          En este caso, el ejercicio ya se nos da implementado por lo que solamente explicaremos en el apartado de Conclusiones el porqué de los rasgos elegidos.
        }
      \subsection*{Implementación}
        \begin{lstlisting}[language=Python, caption=Implementación del clonador de comportamiento del pacman con rasgos diseñados por nosotros]
          
      \subsection*{Conclusiones}
        \paragraph*{}{

        }
  \chapter{Resultados}
    \section{Casos de prueba}
      \subsection{Perceptron}
        \begin{lstlisting}[language=Python, caption=Ejecución del perceptron]
python dataClassifier.py -c perceptron
Doing classification
--------------------
data:           digits
classifier:             perceptron
using enhanced features?:       False
training set size:      100
Extracting features...
Training...
Starting iteration  0 ...
Starting iteration  1 ...
Starting iteration  2 ...
Validating...
55 correct out of 100 (55.0%).
Testing...
48 correct out of 100 (48.0%).

python dataClassifier.py -c perceptron -i 9  
Doing classification
--------------------
data:           digits
classifier:             perceptron
using enhanced features?:       False
training set size:      100
Extracting features...
Training...
Starting iteration  0 ...
Starting iteration  1 ...
Starting iteration  2 ...
Starting iteration  3 ...
Starting iteration  4 ...
Starting iteration  5 ...
Starting iteration  6 ...
Starting iteration  7 ...
Starting iteration  8 ...
Validating...
56 correct out of 100 (56.0%).
Testing...
54 correct out of 100 (54.0%).
        \end{lstlisting}
      \subsection{Clonando el Comportamiento del Pacman}
        \begin{lstlisting}[language=Python, caption=Ejecución del clonador de comportamiento del pacman]
python dataClassifier.py -c perceptron -d pacman
Doing classification
--------------------
data:           pacman
classifier:             perceptron
using enhanced features?:       False
training set size:      100
Extracting features...
Training...
Starting iteration  0 ...
Starting iteration  1 ...
Starting iteration  2 ...
Validating...
83 correct out of 100 (83.0%).
Testing...
80 correct out of 100 (80.0%).

python dataClassifier.py -c perceptron -d pacman -i 9
Doing classification
--------------------
data:           pacman
classifier:             perceptron
using enhanced features?:       False
training set size:      100
Extracting features...
Training...
Starting iteration  0 ...
Starting iteration  1 ...
Starting iteration  2 ...
Starting iteration  3 ...
Starting iteration  4 ...
Starting iteration  5 ...
Starting iteration  6 ...
Starting iteration  7 ...
Starting iteration  8 ...
Validating...
83 correct out of 100 (83.0%).
Testing...
80 correct out of 100 (80.0%).
        \end{lstlisting}
      \subsection{Clonando el Comportamiento del Pacman con rasgos diseñados por nosotros}
    \cpsection{Autograder}
        \begin{lstlisting}[language=Python, caption=Ejecución del autograder]
python autograder.py
Extracting features...
Extracting features...
Starting on 11-8 at 18:30:14

Question q2
===========

Starting iteration  0 ...
Starting iteration  1 ...
Starting iteration  2 ...
Starting iteration  3 ...
79 correct out of 100 (79.0%).
*** PASS: test_cases/q2/grade.test (4 of 4 points)
***     79.0 correct (4 of 4 points)
***         Grading scheme:
***          < 70:  0 points
***         >= 70:  4 points

### Question q2: 4/4 ###


Question q3
===========

Starting iteration  0 ...
Starting iteration  1 ...
Starting iteration  2 ...
Starting iteration  3 ...
Starting iteration  4 ...
80 correct out of 100 (80.0%).
*** PASS: test_cases/q3/contest.test (2 of 2 points)
***     80.0 correct (2 of 2 points)
***         Grading scheme:
***          < 70:  0 points
***         >= 70:  2 points
Starting iteration  0 ...
Starting iteration  1 ...
Starting iteration  2 ...
Starting iteration  3 ...
Starting iteration  4 ...
72 correct out of 100 (72.0%).
*** PASS: test_cases/q3/suicide.test (2 of 2 points)
***     72.0 correct (2 of 2 points)
***         Grading scheme:
***          < 70:  0 points
***         >= 70:  2 points

### Question q3: 4/4 ###


Question q4
===========

Starting iteration  0 ...
Starting iteration  1 ...
Starting iteration  2 ...
Starting iteration  3 ...
95 correct out of 100 (95.0%).
*** PASS: test_cases/q4/contest.test (2 of 2 points)
***     95.0 correct (2 of 2 points)
***         Grading scheme:
***          < 90:  0 points
***         >= 90:  2 points
Starting iteration  0 ...
Starting iteration  1 ...
Starting iteration  2 ...
Starting iteration  3 ...
85 correct out of 100 (85.0%).
*** PASS: test_cases/q4/suicide.test (2 of 2 points)
***     85.0 correct (2 of 2 points)
***         Grading scheme:
***          < 80:  0 points
***         >= 80:  2 points

### Question q4: 4/4 ###


Finished at 18:30:29

Provisional grades
==================
Question q2: 4/4
Question q3: 4/4
Question q4: 4/4
------------------
Total: 12/12

Your grades are NOT yet registered.  To register your grades, make sure
to follow your instructor's guidelines to receive credit on your project.           
        \end{lstlisting}
\end{document}